% HD-F1-T30. GPT-6 Astra Pro; /root/transport; 22 November 2026 UTC.
\section[HD-F1-T30: symplectic multiplicities eight and nine]{HD-F1-T30: the exact factor maps exclude multiplicities eight and nine}
\label{sec:f1-symplectic-eighth-ninth}

The scalar-centre symplectic-type normal subgroup in T17
must have isotypic multiplicity at least ten on every
remaining original T2 source. T17 and T20 excluded the
first seven multiplicities. We now exclude eight and nine
without a classification or a finite normalizer table.

Retain the actual normal subgroup \(R\), scalar cyclic
centre \(Z_R\), and nondegenerate alternating commutator
pairing on \(Q=R/Z_R\). The original coefficient fields
make \(Q\) a symplectic vector space over \(\mathbf F_2\).
Use the actual evaluation isomorphism of T17:
\[
 W\otimes M\xrightarrow{\sim}V_{\mathbf C},\qquad
 d=\dim W=2^a,\quad k=\dim M,\quad m=dk,\quad a\ge1.
 \tag{F1T30.1}
\]
The original inertia tuple has product one, generates
the original coefficient group, has \(b\ge12\) nonidentity
entries, and satisfies
\[
 \sum_{j=1}^b r_j=2dk+6,\qquad
 r_j=\operatorname{rank}(g_j-1)\ge5.
 \tag{F1T30.2}
\]
Put \(C=C_G(R)\) and let \(n\) be the number of entries
outside \(C\). T17 gives \(n\ge3\).

\paragraph{1. The two actual factor tuples and their phases.}
The intertwiners constructed in T20 give the actual
identities \(g_j=A_j\otimes T_j\) under (F1T30.1).
Their only permitted scalar change is
\[
             (A_j,T_j)\longmapsto
                         (\lambda^{-1}A_j,\lambda T_j).
 \tag{F1T30.3}
\]
For \(g_j\in C\), keep \(A_j=1\).
For every other entry, choose a largest eigenspace
of \(A_j\), with eigenvalue \(\lambda\), and apply
(F1T30.3). Thus that largest eigenspace has eigenvalue
one. There is no separate normalization of \(T_j\).
The finite-order argument in T20 makes both factors
semisimple; the scalar change preserves this property.
Write
\[
      u_j=\operatorname{rank}(A_j-1),\qquad
      v_j=\operatorname{rank}(T_j-1)
 \tag{F1T30.4}
\]
for these precise choices. Outside \(C\), \(u_j\) is the
minimum eigenspace codimension of \(A_j\), and T20 gives
\(u_j\ge d/4\). Inside \(C\), \(u_j=0\) and
\(r_j=dv_j\), with \(v_j\ge1\).

Both factor tuples act irreducibly. A common invariant
subspace \(M'\subset M\) gives the original invariant
subspace \(W\otimes M'\), and is therefore zero or all
of \(M\). For \(W\), express any \(r\in R\) as a word
in the original generating \(g_j\) and their inverses.
The corresponding word in the factors satisfies
\[
 A_{\mathrm{word}}\otimes T_{\mathrm{word}}
                       =\rho_W(r)\otimes1.
\]
Equality of these invertible simple tensors gives
\(A_{\mathrm{word}}=z\rho_W(r)\) for a scalar \(z\).
A common invariant subspace for the \(A_j\) is therefore
\(R\)-invariant, hence zero or all of the irreducible
\(W\). Finally the original product-one identity gives
\[
             A_1\cdots A_b=z1_W,\qquad
             T_1\cdots T_b=z^{-1}1_M .
 \tag{F1T30.5}
\]
These are identities on the actual factors, with no
assertion that the scalar \(z\) equals one.

\paragraph{2. The strict rank bound for an irreducible product-scalar tuple.}
For any irreducible tuple \(B_1,\ldots,B_s\) on a
complex vector space of dimension \(h>1\), with
\(\prod B_i=z1\), one has
\[
                \sum_i\operatorname{rank}(B_i-1)\ge h+1.
 \tag{F1T30.6}
\]
Here is the complete proof. If \(z=1\), put
\(P_i=B_1\cdots B_i\), \(P_0=1\), and use the maps
\[
 \begin{aligned}
 F:U&\longrightarrow\bigoplus_i\operatorname{Im}(B_i-1),
       &x&\longmapsto((B_i-1)x)_i,\\
 D:\bigoplus_i\operatorname{Im}(B_i-1)&\longrightarrow U,
       &(x_i)_i&\longmapsto\sum_iP_{i-1}x_i .
 \end{aligned}
\]
Irreducibility makes \(F\) injective and \(D\) surjective:
the latter's annihilator is exactly the common invariant
covectors, as follows successively from
\(\phi P_{i-1}(B_i-1)=0\).
Telescoping gives \(DF=P_s-1=0\). Thus the sum of ranks
is at least \(2h\), which is at least \(h+1\).

For \(z\ne1\), telescoping already gives a sum of ranks
at least \(h\). Suppose equality holds. Set
\[
 C_i=P_{i-1}(B_i-1),\qquad U_i=\operatorname{Im}C_i.
\]
The equality \(\sum_i C_i=(z-1)1\) and the sum of their
ranks \(h\) show that \(U=\bigoplus_iU_i\).
Let \(\pi_i\) be the projections for this actual direct
sum. Comparing the components of that equality gives
\[
 C_i=(z-1)\pi_i,\qquad
 P_j=1+(z-1)\sum_{i\le j}\pi_i.
\]
Therefore \(B_j=P_{j-1}^{-1}P_j\) acts as \(z\) on
\(U_j\) and as one on all the other summands.
With at least two nonzero summands the tuple is
reducible. With one nonzero summand every \(B_j\) is
scalar, also contradicting irreducibility when \(h>1\).
This excludes equality and proves (F1T30.6).
Applied to the two actual tuples in (F1T30.5), it gives
\[
                  \sum_j u_j\ge d+1,\qquad
                  \sum_j v_j\ge k+1.
 \tag{F1T30.7}
\]

\paragraph{3. An exact gap in the symplectic intertwiner's eigenspaces.}
Every nonscalar intertwiner \(A\) arising here has
minimum eigenspace codimension
\[
                       u=d/4\quad\hbox{or}\quad u\ge3d/8.
 \tag{F1T30.8}
\]
We prove this through the actual conjugation maps.
Choose an \(R\)-invariant positive Hermitian form on
\(W\). Since \(A\) intertwines unitary \(R\)-matrices,
\(A^\dagger A\) commutes with \(R\) and is a positive
scalar. After the normalization of paragraph1, \(A\)
is unitary.

Choose representatives \(r_q\in R\), \(q\in Q\).
The \(d^2\) operators \(\rho_W(r_q)\) are an orthogonal
basis of \(\operatorname{End}(W)\). Indeed distinct
representatives have noncentral quotient, and its trace
is zero by the actual commutator character argument of
T17; each operator has squared Hermitian norm \(d\).
Conjugation by \(A\) permutes these operators up to
unit scalars, by the induced symplectic automorphism
\(h\) of \(Q\).
If \(r=\operatorname{rank}_{\mathbf F_2}(h-1)\), the
trace of this actual monomial action therefore gives
\[
 |\operatorname{Tr}A|^2
   =\operatorname{Tr}(\operatorname{Ad}A)
   \le |\ker(h-1)|=\frac{d^2}{2^r}.
 \tag{F1T30.9}
\]
The left side is real and nonnegative, so the last
inequality also follows by taking the absolute value
of the sum of the unit contributions on fixed labels.

Suppose \(u<3d/8\). Its normalized eigenvalue one
has multiplicity \(d-u\), so the triangle inequality
gives \(|\operatorname{Tr}A|\ge d-2u>d/4\).
Equation (F1T30.9) implies \(r\le3\).
Let \(F=\ker(h-1)\). Symplectic orthogonality gives
\(F^\perp=\operatorname{Im}(h-1)\), so the radical of
\(F\) has dimension at most \(r\). Symplectic
Gram--Schmidt therefore supplies an actual nondegenerate
subspace \(E\subset F\) with
\[
              \dim E=2(a-t),\qquad 0\le t\le r\le3.
 \tag{F1T30.10}
\]
Take its full preimage \(R_E\subset R\).
The same averaged-Hom argument as in T17 gives the
actual restriction evaluation
\[
                   W_E\otimes N\xrightarrow{\sim}W,
             \qquad \dim W_E=2^{a-t},\quad \dim N=2^t.
 \tag{F1T30.11}
\]
For every \(r_0\in R_E\), the commutator
\([A,\rho_W(r_0)]\) is a scalar in \(Z_R[2]\), because
\(h\) fixes \(E\) pointwise. These scalars form a
character of \(E\). The perfect pairing on \(E\)
gives \(w\in R_E\) with that same character under
inner conjugation. Consequently, under (F1T30.11),
\[
                         A=\rho_{W_E}(w)\otimes B
 \tag{F1T30.12}
\]
for an actual invertible \(B\).
If \(w\) is noncentral in \(R_E\), its two eigenspaces
on \(W_E\) both have half the dimension. Equation
(F1T30.12) would bound every eigenspace of \(A\) by
\(d/2\), contrary to \(d-u>5d/8\).
Thus \(w\) is central, and every eigenspace dimension
of \(A\) is a multiple of \(2^{a-t}\).
It follows that \(u\) is a multiple of \(d/8\).
The already proved \(u\ge d/4\) leaves only \(u=d/4\)
in the interval under consideration. This proves
(F1T30.8), including the case \(E=0\).

\paragraph{4. Exact tensor fixed-space costs.}
For every outside entry the actual tensor eigenvalue
decomposition gives \(r_j\ge ku_j\): each eigenspace
of \(A_j\) has dimension at most \(d-u_j\).
Together with (F1T30.7), this yields
\[
                    \sum_{g_j\notin C}r_j\ge k(d+1).
 \tag{F1T30.13}
\]
We shall also use a bound retaining the phase of \(T_j\):
\[
           r_j\ge ku_j+(d-2u_j)v_j.
 \tag{F1T30.14}
\]
Indeed its eigenvalue-one space has dimension \(k-v_j\)
and pairs with the specified \((d-u_j)\)-dimensional
eigenvalue-one space of \(A_j\).
All other eigenspaces of \(A_j\) together have dimension
\(u_j\), so their contribution is at most \(u_jv_j\).
Thus
\(\dim V^{g_j}\le(d-u_j)(k-v_j)+u_jv_j\),
which is precisely (F1T30.14).
This proof never changes the original tensor's fixed
eigenvalue to an independently chosen scalar.

\paragraph{5. Multiplicity eight.}
For \(d=2\) one has \(m=16<27\).
For \(d=4\), inside-centralizer ranks are multiples
of four and at least five, hence at least eight;
outside ranks are at least \(kd/4=8\).
Twelve entries then cost at least \(96>70\).
It remains that \(d\ge8\).
The elementary inside and outside costs give
\[
                   16d+6\ge d(b+n).
 \tag{F1T30.15}
\]
Since \(b\ge12\), \(n\ge3\), and \(d\ge8\), one has
\(b+n\le16\). For \(n=3\), (F1T30.13) and the at
least nine inside entries give a cost at least
\(8(d+1)+9d>16d+6\).
For \(n=4\), necessarily \(b=12\), and the cost is
at least \(8(d+1)+8d=16d+8>16d+6\).
These exhaust (F1T30.15), excluding multiplicity eight.

\paragraph{6. Multiplicity nine: the complete integer possibilities.}
Again \(d=2\) has \(m<27\). For \(d=4\), inside
ranks are at least eight and outside ranks at least
nine, so already \(12\cdot8>78\).
Suppose \(d\ge8\).
The elementary costs now give
\[
                 18d+6\ge d\left(b+\frac54n\right).
 \tag{F1T30.16}
\]
If \(n=3\), (F1T30.13) gives the contradiction
\(18d+6\ge9(d+1)+9d=18d+9\).
If \(n=4,b\ge13\), the same contradiction holds.
If \(n\ge6\), the right side of (F1T30.16) is at
least \(19.5d>18d+6\).
For \(n=5,b\ge13\), it is at least
\(19.25d>18d+6\).
Thus the only cases left to examine are
\[
                           b=12,\qquad n=4\ \hbox{or}\ 5.
 \tag{F1T30.17}
\]

For \(n=5\), (F1T30.16) gives \(d\le24\), hence
\(d=8\) or \(16\). The inequalities
\[
 \frac{5d}{4}\le\sum_{\mathrm{outside}}u_j
                \le\left\lfloor\frac{11d+6}{9}\right\rfloor
\]
force all five \(u_j=d/4\) in both cases.
Writing \(V_{\rm in}=\sum_{\mathrm{inside}}v_j\)
and \(V_{\rm out}=\sum_{\mathrm{outside}}v_j\),
(F1T30.14) gives
\[
 d(V_{\rm in}-7)+\frac d2V_{\rm out}
                                      \le6-\frac d4.
 \tag{F1T30.18}
\]
Since \(V_{\rm in}\ge7\), it follows that
\(V_{\rm in}=7\) and \(V_{\rm out}\le1\).
This contradicts \(\sum v_j\ge k+1=10\).

For \(n=4\), the strict factor inequality
\(\sum u_j\ge d+1\) and (F1T30.8) imply
\(\sum u_j\ge9d/8\). Hence
\[
             18d+6\ge9\frac{9d}{8}+8d
                         =\frac{145d}{8},\qquad d\le48.
 \tag{F1T30.19}
\]
Thus \(d=8,16,32\). In these three cases,
\[
 \left\lfloor\frac{10d+6}{9}\right\rfloor=\frac{9d}{8}.
\]
Consequently the four \(u_j\), in some order, are
\[
                        (3d/8,d/4,d/4,d/4).
 \tag{F1T30.20}
\]
Equation (F1T30.14) now gives
\[
 d(V_{\rm in}-8)+\frac d4 V_{\rm out}
                                      \le6-\frac d8 .
 \tag{F1T30.21}
\]
For \(d=16\) or \(32\) this forces
\(V_{\rm in}=8\), \(V_{\rm out}\le1\), again
contradicting their sum at least ten.

\paragraph{7. The final dimension \(d=8,k=9\), with the actual phases.}
Here (F1T30.21) forces \(V_{\rm in}=8\) and
\(V_{\rm out}=2\).
The exact coefficient in (F1T30.14) is two for
the exceptional \(u_j=3\), and four for each of the
other three entries. The available excess is five.
Thus both outside \(v_j\)-dimensions occur at the
exceptional entry; the other three have \(v_j=0\).
All eight inside entries have \(v_j=1\).

The three ordinary outside tensors have rank
\(9\cdot2=18\) each, and the inside sum is \(8\cdot8=64\).
The total \(2\cdot72+6=150\) forces the exceptional
tensor to have rank \(32\), hence fixed dimension \(40\).
In its actual chosen phase, \(A\) has eigenvalue one
with multiplicity five and \(T\) has eigenvalue one
with multiplicity seven. Their contribution is \(35\).
The remaining three \(A\)-dimensions and two
\(T\)-dimensions would have to contribute five.
If those two \(T\)-dimensions form one eigenspace,
their contribution is even. If they form two
distinct eigenspaces, their contribution is at most
the sum of the remaining \(A\)-dimensions, namely three.
Neither gives five.

This matching is also forced without the normalization
notation: if the largest spaces of dimensions five
and seven did not pair as inverse eigenvalues, the
first could pair with at most two \(T\)-dimensions,
and the other three \(A\)-dimensions with at most
seven each. The fixed dimension would be at most
\(5\cdot2+3\cdot7=31<40\).
Thus the contradiction keeps the actual tensor
eigenvalue one in every formulation.
This excludes the last multiplicity-nine case.

\paragraph{Original source maps and precise scope.}
The only additional linear maps are the actual
factor intertwiners with inverse scalar changes,
the telescoping component projections, conjugation
on the explicit operator basis of \(R/Z_R\), and
restriction evaluation on the nondegenerate subspace
\(E\). Both factor tuples are used solely inside
this proved matrix contradiction. No factor is
substituted for a curve coefficient or a detector.
The original connected cover, coefficient local
system, rational primitive corner and restricted
theta polarization stay fixed. Their parabolic
cup pairing, Poincare biextension, period derivative,
cycle-class and Grothendieck/Adams realizations
carry the identity maps with multiplier one.
The assertion applies to the exact scalar-centre
normal-subgroup hypotheses of T17; other normal
subgroup structures are not claimed to be classified.
Together with T17 and T20 it proves the effective
remaining multiplicity bound \(k\ge10\).
